\documentclass[11pt]{article}

\usepackage[margin=1in]{geometry}
\usepackage{amsmath,amssymb,mathtools}
\usepackage{booktabs}
\usepackage{array}
\usepackage{enumitem}
\usepackage{hyperref}
\usepackage{microtype}

\hypersetup{
  colorlinks=true, linkcolor=blue, citecolor=blue, urlcolor=blue,
  pdftitle={Invariant-Driven Recursive Refinement for Breeding-Blanket Design},
  pdfauthor={Michael A. Russell}
}

\title{
\textbf{Invariant-Driven Recursive Refinement for Breeding-Blanket Design}\\[0.4em]
\large A Preregistered, Integer-Exact Reproducibility Study on a Reduced-Order Surrogate\\[0.2em]
\normalsize Revised manuscript, revision 5 (supersedes earlier versions)
}

\author{Michael A. Russell}
\date{October 29, 2026 (revision 5)}

\begin{document}
\maketitle

\begin{abstract}
We apply Russell Recursive Refinement (R$^3$) to the coupled design of a water-cooled
liquid-breeder fusion blanket slice, trading tritium breeding ratio (TBR) against coolant
pressure drop under eleven neutronic, thermal, structural, shielding, manufacturing, and
inventory constraints. The investigation is implemented as a reproducibility capsule in
which every quantity is an exact WAD-18 integer (scale $10^{18}$, floor rounding); an
automated audit finds no floating-point use in 35 source files or 89 output files.
Five protocol versions were each frozen before their first execution and are retained in full.

Within the reduced-order surrogate model, a margin-aware search (protocol v1.1.0) identifies a design
that raises TBR from 1.160 to 1.302 and lowers pressure drop from 20.8 to 17.6~kPa
relative to a fixed reference, meets every declared limit with at least 5.39\% margin,
and remains feasible and superior to the reference at all 64 corners of a six-parameter
uncertainty box. The first protocol (v1.0.1) produced no admissible design: every
Pareto-optimal design lay on a constraint boundary, and the safety margin had been applied
only after the search.

The original R$^3$ did not outperform NSGA-II. We therefore refined R$^3$ in three
confirmatory cycles. Each cycle diagnosed the previous failure on development data, and each
used a fresh held-out set that was declared before development began; the decision rule never
changed. Cycle 1 scored 2 of 8 held-out cells (FAIL) and cycle 2 scored 5 of 8 (FAIL). Cycle 3,
whose R$^3$ escalates to coupled two-variable moves whenever an iteration stalls, scored 7 of 8
(PASS). Because three attempts were made, an idealised null gives about a 10\% chance of such a
pass, so the evidence for R$^3$'s advantage is moderate rather than decisive. Physical relevance is
not established: the model is a surrogate with unvalidated constants. Within it, proof
obligations 1--5 now pass. An integer multigroup transport model using verified FENDL-3.2c data confirms the
direction of the design improvement. A 1-D slab leaves the absolute TBR bracketed by the plasma-side
boundary treatment, but an integer cylindrical model that includes the plasma chamber resolves it: TBR is 1.268
for the reference and 1.485 for the selected design, with 3.3\% numerical uncertainty. All three designs meet
the requirement, and the surrogate is conservative by 9--13\%. The work is not a demonstrated fusion-engineering
advance: 3-D effects enter only through an assumed coverage factor, and no experimental benchmark has been compared.
\end{abstract}

\noindent\textbf{Keywords:} breeding blanket; tritium breeding ratio; multi-objective optimization;
preregistration; held-out evaluation; integer-exact arithmetic; reproducibility; recursive refinement.

\section{Purpose and Summary of Findings}

The original manuscript set out to advance fusion engineering by applying R$^3$ to integrated
breeding-blanket performance and to support the result with an executable capsule. It defined
six proof obligations. This revision reports what the executed capsule established for each
of them; the verdicts were fixed by rules declared before each execution.

\begin{table}[h]
\centering
\caption{Proof-obligation outcomes by protocol version.}
\label{tab:po}
\begin{tabular}{lccccc}
\toprule
Obligation & v1.0.1 & v1.1.0 & v1.2.0 & v1.3.0 & v1.4.0 \\
\midrule
PO1 reproducible execution & PASS & PASS & PASS & PASS & PASS \\
PO2 constraint satisfaction & FAIL$^a$ & PASS & PASS & PASS & PASS \\
PO3 engineering improvement & FAIL$^a$ & PASS & PASS & PASS & PASS \\
PO4 contribution of R$^3$ & FAIL & FAIL & FAIL 2/8 & FAIL 5/8 & PASS 7/8 \\
PO5 robustness (64 corners) & FAIL$^a$ & PASS & PASS & PASS & PASS \\
PO6 physical relevance & n.e. & n.e. & n.e. & n.e. & n.e. \\
\bottomrule
\multicolumn{6}{l}{\footnotesize $^a$ No design met the selection rule. PO4 in v1.2.0--v1.4.0 is decided on held-out cells (6 of 8 required).}\\
\multicolumn{6}{l}{\footnotesize n.e.: not established.}
\end{tabular}
\end{table}

PO1--PO5 hold in v1.4.0, with PO2, PO3 and PO5 holding \emph{within the surrogate model}. PO4 passed
on the third confirmatory cycle, a qualification that Section~\ref{sec:mult} quantifies. PO6 cannot be
satisfied by this capsule. These outcomes are not grounds to describe the work as a
fusion-engineering breakthrough, and this revision does not do so.

\section{The Engineering Problem}

The breeding blanket breeds tritium, captures neutron energy, removes heat, shields the
magnets, and must remain structurally sound and manufacturable. These functions compete for
the same radial space and volume fractions. The design vector is
\[
d = (t_b,\ e,\ t_m,\ f_s,\ f_c,\ D),
\]
comprising breeding-zone thickness, Li-6 enrichment, multiplier thickness, stiffener fraction,
coolant flow-area fraction, and tube diameter, each on a declared grid
($31\times13\times9\times21\times17\times17\approx2.2\times10^7$ designs). The objectives are
\[
\max_{d\in\mathcal F} T(d), \qquad \min_{d\in\mathcal F}\Delta p(d),
\]
and $\mathcal F$ is defined by eleven constraints: $T\ge1.10$, $\Delta p\le200$~kPa,
steel and breeder temperatures $\le550\,^\circ$C, tube hoop stress $\le100$~MPa, magnet
fast-flux attenuation $\ge e^{-5.5}$, shield thickness $\ge0.10$~m, stiffener fraction above a
depth-dependent minimum, tube ligament $\ge3$~mm, multiplier tritium inventory $\le3$~g/m$^2$,
and coolant velocity $\le3$~m/s.

\subsection{The model and its status}

The model operator $y=\mathcal M(d,\theta,u)$ is a reduced-order surrogate built from
one-dimensional saturation forms for breeding and energy capture, exponential attenuation for
shielding, the Blasius friction and Dittus--Boelter heat-transfer correlations, and closed-form
conduction, hoop-stress, and geometry relations. Its constants are declared, plausible
magnitudes, not values from qualified datasets. It contains no neutron transport, nuclear-data
library, spatial power profile, or transient analysis. Every result below is a statement about
this model.

\section{Method}

\subsection{R$^3$ as implemented}

R$^3=I\circ V\circ F\circ A$ is realised as a constraint-aware, Pareto-archive pattern search.
\begin{itemize}[leftmargin=1.2em]
  \item \textbf{Analyze:} select focus designs (front extremes and the largest hypervolume gap),
  record binding constraints, and rank move directions by their recorded success.
  \item \textbf{Refine:} propose coordinate moves at the current scale, extrapolate each design's
  last successful move, and attempt repair combinations for promising designs that violate
  exactly one constraint.
  \item \textbf{Validate:} evaluate every constraint, admit only feasible non-dominated designs, and
  check the archive invariants, including that hypervolume never decreases.
  \item \textbf{Iterate:} keep the scale on a hypervolume gain and halve it otherwise. A fixed point
  $\Phi(A)=A$ is reached at unit scale with no gain.
\end{itemize}
Version 1.1.0 added a multiscale restart rule. Version 1.2.0 added Latin-hypercube seeding~\cite{McKay1979}
and re-seeding at local fixed points. Version 1.3.0 added \emph{neighbourhood escalation}: when the
single-variable neighbourhood of the front is exhausted, R$^3$ proposes coupled two-variable moves and
then recombinations of front designs before any random re-seeding. Version 1.4.0 also escalates after
any iteration that brings no gain.

\subsection{Baselines and comparison}

The baselines are NSGA-II~\cite{Deb2002}, with constrained domination and crowding distance, and
uniform random search. Every method receives the same budget of unique evaluations, and each
starts from the reference design. Fronts are compared by exact two-dimensional hypervolume with reference point
$(T,\Delta p)=(1.10,\,200\text{ kPa})$ and $\Delta p$ normalised to make the measure dimensionless.

\subsection{Integer-exact computation}

All arithmetic uses WAD-18 integers with floor rounding and signed 256-bit result bounds.
The exponential, square root, and fractional powers ($Re^{-1/4}$, $Re^{4/5}$, $Pr^{2/5}$) are
computed exactly to a single final floor. They match $e$, $e^{-1}$, $e^{2}$, $\sqrt2$, and
$2^{1/4}$ to the last of 18 digits, and integer roots satisfy $r^k\le n<(r+1)^k$ by construction.
An abstract-syntax audit rejects float literals, true division, non-integer exponents,
approximate-numeric modules, and float tokens in every output file.

\subsection{Preregistration and the evidence chain}

Each protocol version was written and hashed before its first execution. Every change made
after an execution is recorded with its effect on the results, and earlier studies are
re-executed and byte-compared on every build. These freezes are self-attested within one
working session; no external timestamp authority was used. Each study writes its complete
evaluation logs, iteration records, uncertainty corners, sensitivity table, output-level
checks, and proof-obligation verdicts, and a manifest records the SHA-256 of every file.

\section{Results}

\subsection{v1.0.1: optimal designs sit on constraints}

R$^3$ reached a unit-scale fixed point after 874 of 1000 evaluations. Sixteen of its twenty
front designs dominated the reference, with TBR up to about 1.36 at lower pressure drop.
However, the predeclared selection rule required every normalised constraint slack to be at
least 0.05, and none of these designs met it: each lay on or near the shield-attenuation,
stiffener, multiplier-inventory, or hoop-stress limit. This is the expected behaviour of an
unconstrained Pareto search, which drives designs onto the boundaries of the feasible set.
Because the margin was applied only after the search, no admissible candidate existed.

\subsection{v1.1.0: margin-aware search}

Version 1.1.0 moved the 5\% margin into the feasibility test used during search, for every method.
The selection rule then returned the design in Table~\ref{tab:cand}.

\begin{table}[h]
\centering
\caption{Selected design versus reference (nominal $\theta$, surrogate model).}
\label{tab:cand}
\begin{tabular}{lrr}
\toprule
Quantity & Reference & Selected \\
\midrule
TBR & 1.160 & 1.302 \\
Pressure drop (kPa) & 20.80 & 17.56 \\
Peak steel temperature ($^\circ$C) & 412.9 & 422.2 \\
Peak breeder temperature ($^\circ$C) & 497.1 & 519.2 \\
Tube hoop stress (MPa) & 77.5 & 85.3 \\
Attenuation exponent & 5.889 & 5.797 \\
\midrule
Design $(t_b, e, t_m, f_s, f_c, D)$ & \multicolumn{2}{r}{$(0.47\,\text{m}, 0.90, 0.05\,\text{m}, 0.105, 0.030, 11\,\text{mm})$} \\
\bottomrule
\end{tabular}
\end{table}

The smallest constraint slack is 5.39\% (shield attenuation). Across all 64 corners of the
uncertainty box, which perturbs breeding potential by $\pm3\%$, attenuation length by $\pm5\%$,
multiplier gain by $\pm10\%$, friction by $\pm10\%$, heat transfer by $\pm15\%$, and breeder
conductivity by $\pm10\%$, the design stays feasible and dominates the reference. The robustness is
narrow, however: in the worst corner the breeder-temperature slack falls to 0.47\%. Most of the
TBR gain comes from the 5~cm multiplier layer, which the surrogate represents generically.

\subsection{Search performance, cycle 1: the original R$^3$ does not outperform NSGA-II}

\begin{table}[h]
\centering
\caption{Hypervolume, development instance, budget 1000 (NSGA-II and random: seeds 1--5).}
\label{tab:hv}
\begin{tabular}{lccc}
\toprule
Method & v1.0.1 & v1.1.0 & v1.2.0 \\
\midrule
R$^3$ & 0.2642 & 0.2015 & 0.2114$^c$ \\
NSGA-II median (min--max) & 0.2665 (0.2495--0.2710) & 0.2082 (0.1960--0.2141) & as v1.1.0 \\
Random median & 0.1982 & 0.1299 & as v1.1.0 \\
\bottomrule
\multicolumn{4}{l}{\footnotesize $^c$ Single run, seed 0, on the development instance. This does not decide PO4.}
\end{tabular}
\end{table}

For v1.2.0, R$^3$ itself was refined using R$^3$'s own cycle. A logged development harness
evaluated 21 variants on the v1.1.0 conditions. One finding came early: a single-run
``win'' (variant~7) vanished when the seeding seed changed, so later variants were compared
as distributions over R$^3$'s own seeds. The chosen design, 20 Latin-hypercube points plus
re-seeding at local fixed points, reached a development median of 0.2104 against the NSGA-II
median of 0.2082. That margin is small, and it was selected from 21 attempts.

The design was then frozen, together with a held-out test it had never seen: four instances
(wall load $\pm20\%$, radial build $+5$~cm, breeder conductivity $-20\%$), budgets of 1000 and 2000,
and seven new seeds per method. Table~\ref{tab:ho} gives the outcome.

\begin{table}[h]
\centering
\caption{Cycle-1 held-out test (v1.2.0): median hypervolume over seeds 101--107.}
\label{tab:ho}
\begin{tabular}{lrcccc}
\toprule
Instance & Budget & R$^3$ & NSGA-II & Random & R$^3$ ahead \\
\midrule
H1 high wall load & 1000 & 0.1298 & 0.2013 & 0.1164 & no \\
H1 high wall load & 2000 & 0.1818 & 0.2039 & 0.1343 & no \\
H2 low wall load & 1000 & 0.2009 & 0.2164 & 0.1459 & no \\
H2 low wall load & 2000 & 0.2009 & 0.2194 & 0.1541 & no \\
H3 wider radial build & 1000 & 0.2384 & 0.2433 & 0.1780 & no \\
H3 wider radial build & 2000 & 0.2494 & 0.2493 & 0.1910 & yes \\
H4 low breeder conductivity & 1000 & 0.1475 & 0.2022 & 0.1246 & no \\
H4 low breeder conductivity & 2000 & 0.2084 & 0.2063 & 0.1413 & yes \\
\bottomrule
\end{tabular}
\end{table}

R$^3$ led NSGA-II in 2 of 8 cells, against the 6 the rule required, so cycle 1 failed. R$^3$ led
random search in all 8 cells. Its median in H2 is unchanged between budgets of 1000 and 2000,
which indicates stagnation rather than slow convergence.

\subsection{Cycles 2 and 3: diagnosing and removing the stagnation}

The cycle-1 cells then became development data. Their iteration logs showed the cause of the
stagnation: after roughly 900 evaluations, every one-variable neighbour of the front had already
been evaluated. R$^3$ made no new evaluations and fell back on random re-seeding. The variables are
coupled (breeder and multiplier compete for radial space; coolant fraction and tube diameter interact),
so the front was optimal only against one-variable moves.

Neighbourhood escalation (v1.3.0) addressed this. A competing hypothesis, refining the least-violating
infeasible designs, was tested and refuted (0 of 9 development cells); it is retained in the log.
Cycle 2 scored 5 of 8 on a new held-out set (H5--H8: lower coverage, a longer coolant path,
pessimistic nuclear parameters, and a hotter, lower-conductivity case; budgets 1500 and 3000;
seeds 201--207). One of its losses was an exact tie, which the unchanged rule counts against R$^3$.

The cycle-2 cells again became development data. Their logs showed long plateaus during which
untested neighbours existed but none improved the front, so escalation fired late. Version 1.4.0 escalates after
any non-improving iteration. Its held-out set (H9--H12) was declared before cycle-3 development began.
The outcome is given in Table~\ref{tab:ho3}.

\begin{table}[h]
\centering
\caption{Cycle-3 held-out test (v1.4.0): median hypervolume over seeds 301--307.}
\label{tab:ho3}
\begin{tabular}{lrcccc}
\toprule
Instance & Budget & R$^3$ & NSGA-II & Random & R$^3$ ahead \\
\midrule
H9 enrichment cost $+40\%$ & 1200 & 0.2059 & 0.2039 & 0.1234 & yes \\
H9 enrichment cost $+40\%$ & 2500 & 0.2069 & 0.2069 & 0.1445 & tie (no) \\
H10 thinner manifold & 1200 & 0.2353 & 0.2268 & 0.1653 & yes \\
H10 thinner manifold & 2500 & 0.2380 & 0.2350 & 0.1730 & yes \\
H11 weaker heat transfer & 1200 & 0.2089 & 0.2062 & 0.1342 & yes \\
H11 weaker heat transfer & 2500 & 0.2132 & 0.2092 & 0.1533 & yes \\
H12 lower load, shorter $\lambda_0$ & 1200 & 0.2367 & 0.2245 & 0.1568 & yes \\
H12 lower load, shorter $\lambda_0$ & 2500 & 0.2383 & 0.2291 & 0.1757 & yes \\
\bottomrule
\end{tabular}
\end{table}

R$^3$ led in 7 of 8 cells, so PO4 passes. In 3 cells the R$^3$ median exceeded the \emph{best}
NSGA-II seed, and in 3 cells all seven R$^3$ seeds exceeded the NSGA-II median.

\subsection{How strong is the PO4 pass?}
\label{sec:mult}

The pass came on the third attempt. Suppose each cell were a fair coin with cells independent, an
idealisation since cells share instances. A single cycle would then reach at least 7 of 8 with
probability $9/256$, and at least one of three cycles would do so with probability
$1-(247/256)^3 = 1707993/16777216 \approx 0.102$. At the seed level, 41 of 56 cycle-3 R$^3$ seeds
beat the NSGA-II median, a tail probability of about $3.4\times10^{-4}$, but seeds within a cell are not
independent, so this overstates the evidence. We read the result as moderate evidence that
escalating R$^3$ is competitive with or better than NSGA-II on this problem family. It does not
show general superiority.

\subsection{The design selected by v1.4.0}

Run on the development instance, the v1.4.0 R$^3$ selects $(0.46\,\text{m}, 0.90, 0.05\,\text{m}, 0.105, 0.030, 11\,\text{mm})$,
which differs from the v1.1.0 design only in breeder thickness. It gives TBR 1.295, $\Delta p$ 18.19~kPa and a
minimum nominal slack of 5.19\%, and it remains feasible and dominating at all 64 corners. Its worst-corner
breeder-temperature slack, however, is only 0.0072\%. On robustness grounds the v1.1.0 design
(worst-corner slack 0.47\%, Table~\ref{tab:cand}) is preferable. A better optimizer does not by itself
yield a more robust design while the selection rule considers only nominal margins.

\section{Transport Check with Evaluated Nuclear Data}
\label{sec:transport}

\paragraph{Data.} FENDL-3.2c multigroup libraries (NJOY2016, 211 neutron groups, 293.6~K) for all 21 nuclides
of the model (H-1, O-16, Li-6/7, Be-9, Cr-50/52/53/54, Fe-54/56/57/58, W-182/183/184/186 and
Pb-204/206/207/208) are read by integer-exact readers. Every library passes four checks:
\begin{itemize}[leftmargin=1.2em]
  \item exact word counts;
  \item elastic transfer out of each group equal to the elastic cross section within 0.1\% below 20~MeV, which
  also fixes the storage layout empirically;
  \item a total cross section equal to the sum of its neutron-emitting and pure-absorption reactions within
  10~ppm;
  \item group values within the weight-independent bounds of the corresponding pointwise ENDF-6 evaluations
  (outside resonance ranges and above 100~eV).
\end{itemize}
The Legendre normalisation and (n,2n) multiplicity were verified from the data rather than assumed.

\paragraph{Solver.} A 1-D slab S$_8$ step-characteristic solver with P$_3$ scattering, written entirely in
WAD-18 integers, solves the 211-group problem with a 14.1~MeV source on the plasma side.
\begin{itemize}[leftmargin=1.2em]
  \item \emph{Verification.} It reproduces analytic pure-absorber transmission to $10^{-18}$, and it conserves
  neutrons to $10^{-10}$ or better in every run.
  \item \emph{Convergence.} At the chosen resolution, angular and Legendre truncation contribute below 0.15\%.
  Halving the mesh resolution changes TBR by 1.4\%, which we attach as the numerical uncertainty.
\end{itemize}
Assumptions the surrogate never stated are declared in a separate specification: material densities, a
beryllium multiplier, and the manifold and shield compositions.

\paragraph{Boundary condition.} The declared vacuum plasma-side boundary gave a reference TBR of 0.45. In that
run, 109\% of the incident current left through the front face: Pb(n,2n) multiplication combined with slow
moderation in lead. In a torus those neutrons re-enter the blanket elsewhere. We therefore added a white
(isotropic re-entry) boundary \emph{after} seeing this result, and we report the two as bounds.

\paragraph{Pre-declared questions.} Four questions were fixed before any non-reference design was computed.
Table~\ref{tab:tr} gives the outcome.

\begin{table}[h]
\centering
\caption{TBR (including coverage 0.82): surrogate versus transport bounds.}
\label{tab:tr}
\begin{tabular}{lccc}
\toprule
Design & Surrogate & Transport, vacuum & Transport, white \\
\midrule
Reference & 1.160 & 0.453 & 1.300 \\
v1.1.0 selected & 1.302 & 0.575 & 1.522 \\
v1.4.0 selected & 1.295 & 0.574 & 1.520 \\
\bottomrule
\end{tabular}
\end{table}

\begin{itemize}[leftmargin=1.2em]
  \item \textbf{Direction (AGREE).} Transport confirms that both selected designs breed more than the reference,
  by 17\% (white) to 27\% (vacuum), far beyond the numerical uncertainty.
  \item \textbf{Magnitude (NO).} The surrogate's gain of about 12\% lies below that range: it undervalues the
  beryllium layer.
  \item \textbf{Requirement (DEPENDS ON BOUNDARY).} Whether TBR $\ge 1.155$ holds depends on the boundary
  treatment.
  \item \textbf{Calibration.} Every surrogate value lies inside the transport bounds.
\end{itemize}

\paragraph{Cylindrical model.} The slab cannot represent the plasma chamber, so we built a 1-D cylindrical
discrete-ordinates solver, again entirely in WAD-18 integers.
\begin{itemize}[leftmargin=1.2em]
  \item \emph{Geometry.} A plasma region of radius 2.9~m carries a uniform 14.1~MeV source. The first wall is
  at 3.1~m, followed by the same layers and materials as before. Neutrons leaving the blanket front cross the
  chamber and re-enter on the far side.
  \item \emph{Numerics.} Step differencing in space and angle keeps the scheme positive and exactly conservative.
  \item \emph{Verification.} The solver reproduces an exactly flat flux in an infinite medium (12 digits, zero
  current), and it conserves neutrons to $10^{-13}$.
\end{itemize}
Its numerical uncertainty is $U = 3.3\%$, dominated by the mesh. Questions and rules were again fixed before
any cylindrical blanket run (Table~\ref{tab:cyl}).

\begin{table}[h]
\centering
\caption{TBR (including coverage 0.82) in the cylindrical model, with the slab bracket.}
\label{tab:cyl}
\begin{tabular}{lccc}
\toprule
Design & Surrogate & Slab bracket & Cylinder \\
\midrule
Reference & 1.160 & 0.453--1.300 & 1.268 \\
v1.1.0 selected & 1.302 & 0.575--1.522 & 1.485 \\
v1.4.0 selected & 1.295 & 0.574--1.520 & 1.483 \\
\bottomrule
\end{tabular}
\end{table}

\begin{itemize}[leftmargin=1.2em]
  \item \textbf{Bracket.} Every cylinder value lies inside its slab bracket.
  \item \textbf{Direction (AGREE).} The selected designs gain 17\%, about five times $U$.
  \item \textbf{Requirement (MET).} All three designs meet TBR $\ge 1.155$ by more than $U$.
  \item \textbf{Surrogate.} It underestimates TBR by 8.5--12.7\%, and its predicted gain falls outside
  $\pm U$ of the computed one.
  \item \textbf{Geometry sensitivity.} Moving the first wall from 2.5~m to 4.0~m changes the reference TBR by
  only $\pm0.1\%$.
\end{itemize}

The design direction found by the search is therefore supported by evaluated physics, and within this model
the requirement is met. What remains open is three-dimensional: the coverage factor (0.82, assumed), the first
wall (absent from every model here), ports and the divertor.

\section{Discussion}

\paragraph{What is established.}
Within the stated surrogate, the capsule establishes three things. First, there are designs that
improve both objectives over the reference while carrying a uniform safety margin and surviving the
declared uncertainty box. Second, every reported number regenerates byte-identically from the frozen
inputs. Third, no floating-point arithmetic is used anywhere in the capsule's own code or outputs.

\paragraph{What the failures teach.}
The v1.0.1 failure shows that a safety margin applied after optimization is the wrong place for it.
Pareto optimality pushes designs onto constraint boundaries, so margins belong inside the search.
The cycle-1 failure shows why a held-out evaluation is necessary: refinements tuned on a single
development problem produced an edge that did not transfer. The improvements that did transfer
(cycles 2 and 3) came from diagnosing a mechanism, neighbourhood exhaustion and stalled plateaus,
rather than from adjusting parameters.

\paragraph{Recursive refinement of the study itself.}
The evidence process was itself refined recursively, through cycles of
analysis, revision, validation, and iteration. Each cycle produced a frozen protocol, an executed
result, and a logged revision. The process stops being informative if its acceptance criteria are
relaxed after the results are seen, so no criterion was changed after its study ran.

\section{Limitations}

\begin{itemize}[leftmargin=1.2em]
  \item \textbf{Surrogate model.} The search uses the surrogate, whose constants are unvalidated.
  Section~\ref{sec:transport} checks only its TBR, in 1-D, and there is no spatial power profile, transient,
  corrosion, or irradiation-embrittlement modelling. Tritium
  recovery enters only through the required TBR, and no plant-level net-electric claim is made.
  \item \textbf{Statistical scale and multiplicity.} Each held-out cycle used seven seeds per method on one problem
  family, and PO4 passed on the third cycle (Section~\ref{sec:mult}). All three held-out sets are now
  used; any further claim requires new held-out conditions.
  \item \textbf{Nominal-margin selection.} The selection rule considers only nominal margins, which allowed a
  design with very small worst-corner slack (v1.4.0). Selecting on worst-corner slack is an obvious next protocol change.
  \item \textbf{Self-attested preregistration.} The freeze records and their hashes are internal to the capsule.
  \item \textbf{Reproduction scope.} Reproduction has been same-machine and self-executed. Independent third-party
  reproduction has not yet taken place.
\end{itemize}

\section{Path to Physical Relevance (PO6)}

The capsule isolates the model operator in one module, so the search, constraint, comparison,
uncertainty, and evidence machinery are unchanged when the model is replaced. Establishing PO6
requires three things. The first is a Monte Carlo transport model of the candidate and reference geometries
(e.g.\ OpenMC, as used for WCLL TBR studies~\cite{Chiletti2026}). The 1-D check of Section~\ref{sec:transport}
and the cylindrical model narrow what remains: the coverage factor, the first wall, and 3-D features. A
Monte Carlo code is one route; it need not use floating point, since a seeded integer generator makes it as
reproducible as the rest of the capsule. Validation still requires comparison with measured tritium production,
for example from blanket mock-up experiments. The second is thermal-hydraulic
analysis at the fidelity reviewed in~\cite{Seo2021}. The third is comparison against benchmark or experimental data
relevant to the materials and conditions. Only then could PO2, PO3, and PO5 be read as statements about a
physical blanket rather than about the surrogate.

\section{Conclusion}

The executed study does not demonstrate a nuclear-fusion engineering advance. It provides a
complete, integer-exact, preregistered and byte-reproducible evidence chain. Within a declared surrogate
model, that chain establishes a margin-carrying, uncertainty-robust design improvement, which integer transport calculations with evaluated nuclear data confirm in
direction and show to meet the TBR requirement in 1-D cylindrical geometry. After three
disclosed confirmatory cycles, it also provides moderate held-out evidence that an escalating R$^3$
outperforms NSGA-II on this problem family. Every failure along the way is recorded rather than
concealed. The next step that could
turn the model-internal result into an engineering one is replacing the surrogate with validated
transport and thermal-hydraulic analysis.

\section*{Data and Code Availability}
The reproducibility capsule (\texttt{fusion\_r3\_capsule}) contains all specifications, source
code, logs, results, freeze records, the development log, and a SHA-256 manifest. It needs only the
Python standard library and is rebuilt with \texttt{./build\_capsule.sh}.

\begin{thebibliography}{99}

\bibitem{DOE2025}
U.S. Department of Energy.
\textit{Fusion Science and Technology Roadmap}. 2025.
\url{https://www.energy.gov/sites/default/files/2025-10/fusion-s%26t-roadmap-101625.pdf}.

\bibitem{Chiletti2026}
M. Chiletti, N. Rebiai, D. Mazed, and F. G\'erardin.
Predictive modelling of tritium breeding ratio (TBR) in EU DEMO WCLL blankets: OpenMC
simulations and analytical insights.
\textit{Fusion Engineering and Design}, 222, 115516, 2026.
doi:10.1016/j.fusengdes.2025.115516.

\bibitem{Seo2021}
S. B. Seo et al.
A review of thermal hydraulics systems analysis for breeding blanket design and future needs
for fusion engineering demonstration facility design and licensing.
\textit{Fusion Engineering and Design}, 172, 112769, 2021.
doi:10.1016/j.fusengdes.2021.112769.

\bibitem{IAEA2025}
International Atomic Energy Agency.
Technical Meeting on Tritium Breeding Blankets and Associated Neutronics. 2025.
\url{https://conferences.iaea.org/event/406/}.

\bibitem{Deb2002}
K. Deb, A. Pratap, S. Agarwal, and T. Meyarivan.
A fast and elitist multiobjective genetic algorithm: NSGA-II.
\textit{IEEE Transactions on Evolutionary Computation}, 6(2):182--197, 2002.

\bibitem{McKay1979}
M. D. McKay, R. J. Beckman, and W. J. Conover.
A comparison of three methods for selecting values of input variables in the analysis of
output from a computer code.
\textit{Technometrics}, 21(2):239--245, 1979.

\end{thebibliography}

\end{document}
